\begin{proof}
Put $u=q^*g$.  Fix $x\in X$, choose $y\in q^{-1}(x)$, and set
$H:=G_y$.  Since $Y$ is Hausdorff, there is an ordinary open neighbourhood
$C$ of $y$ which is $H$-stable and whose translates indexed by $G/H$ are
pairwise disjoint.  The orbit topology and invariant structure sheaf for the
finite analytic quotient identify
\begin{equation}\label{eq:p061-clean-local-quotient-candidate}
                         C/H\simeq q(C),
\end{equation}
where $q(C)$ is an ordinary open neighbourhood of $x$ in $X$.

By Gubler--Jell--Rabinoff Proposition~5.8, after shrinking around $y$ in the
ordinary topology, there are $a_1,\ldots,a_m\in\mathbf R$ and integral PL
functions $h_1,\ldots,h_m$ such that
\begin{equation}\label{eq:p061-local-rpl-decomposition-candidate}
                              u=\sum_{i=1}^m a_i h_i.
\end{equation}
By adjoining the zero integral PL function with coefficient zero if
necessary, we may assume $m\geq1$.
Intersect this ordinary open with $C$ and with all its $H$-translates.  We may
therefore assume it is $H$-stable and contained in $C$.  Choose a strict
affinoid neighbourhood $N_0$ of $y$ in it and put
\[
                          N:=\bigcap_{\gamma\in H}\gamma N_0.
\]
Separatedness makes $N$ a strict affinoid domain.  Since every factor is an
ordinary neighbourhood of the $H$-fixed point $y$, $N$ is an $H$-stable
strict affinoid neighbourhood of $y$.  Write $N^\circ$ for its interior
relative to $C$.

We next record the formal-model construction, including the descent step
because the quotient $N\to N/H$ need not be flat.  Restrict the $h_i$ to
$N$.  By Gubler--Jell--Rabinoff~5.6 and Raynaud cofilteredness, they are model
functions on a common formal model.  This common model may be chosen
$H$-equivariant as follows.  Start with an $H$-stable affine formal model
$\mathfrak N_0$ of $N$.  For a model
$(\mathfrak N_i,\mathfrak L_i)$ inducing $h_i$, Raynaud cofilteredness gives
an admissible formal blowup $\mathfrak P_i\to\mathfrak N_0$ and a morphism
$\mathfrak P_i\to\mathfrak N_i$ inducing the identity on the generic fibre.
Take the schematic closure of the common diagonal generic fibre in the
fibre product
\[
       \prod_{\substack{\mathfrak N_0;\ 1\leq i\leq m,\ \gamma\in H}}
             \gamma^*\mathfrak P_i.
\]
After removing $K^\circ$-torsion, the resulting admissible formal model
$\mathfrak N$ is $H$-equivariant, and projection to the identity factors
pulls the line bundles back to line bundles, still denoted
$\mathfrak L_i\in M(\mathfrak N)$, satisfying
$h_i=h_{\mathfrak L_i}$.  Every $\mathfrak P_i\to\mathfrak N_0$ is
projective.  Thus $\mathfrak N$ is closed in a projective formal scheme over
the affine $\mathfrak N_0$, and its special fibre is quasi-projective.  Hence
Zavyalov Remark~3.3.3 and Theorem~3.3.4 construct the admissible formal
quotient $\pi:\mathfrak N\to\mathfrak X:=\mathfrak N/H$, and Zavyalov
Theorem~4.4.1 identifies the adic generic fibre of $\mathfrak X$ with the
adic quotient of $N$.  Passing to the associated strict Berkovich space, or
equivalently using the invariant affinoid algebra, identifies it with the
Berkovich quotient $N/H$.

For each $i$, define the tensor-induced line bundle
\[
        \mathfrak T_i:=\bigotimes_{\gamma\in H}\gamma^*\mathfrak L_i.
\]
It has a canonical $H$-linearization by reindexing the factors.  We claim
that it descends to a line bundle $\overline{\mathfrak L}_i$ on
$\mathfrak X$.  Indeed, fix a point of $\mathfrak X$.  Its inverse image is
one finite $H$-orbit.  On an affine neighbourhood of that orbit, semilocal
triviality of a rank-one projective module gives an open neighbourhood $V$
of the entire orbit and one frame $s$ of $\mathfrak L_i$ on $V$.  Replacing
$V$ by the intersection of its $H$-translates makes it $H$-stable while
preserving the frame.  Then
\[
                         \bigotimes_{\gamma\in H}\gamma^*s
\]
is an invariant frame of $\mathfrak T_i$.  Its image is an open subset of
$\mathfrak X$, and ratios of any two such invariant frames are invariant
units.  Since the quotient structure sheaf is the invariant structure sheaf,
those ratios descend to units and glue trivial line bundles on
$\mathfrak X$.  This proves the claim.  Tensoring the given generic
trivializations of the $\gamma^*\mathfrak L_i$ gives an invariant analytic
frame of $\mathfrak T_{i,\eta}$.  Its ratios with the invariant formal frames
just constructed are invariant analytic units, and therefore descend and
glue to a trivialization of $\overline{\mathfrak L}_{i,\eta}$.  Thus
$\overline{\mathfrak L}_i$ is indeed a model of the analytically trivial line
bundle, with trivialization pulling back to the tensor-product
trivialization.  Notice that this argument neither uses exactness of invariants nor
divides by $|H|$ in $K^\circ$.

Let $\overline h_i:=h_{\overline{\mathfrak L}_i}$ on $N/H$.  Compatibility
of model functions with pullback gives
\begin{equation}\label{eq:p061-tensor-induction-pullback-candidate}
                  q_N^*\overline h_i
                     =\sum_{\gamma\in H}\gamma^*h_i.
\end{equation}
Consequently the $\mathbf R$-PL function
\[
                       v:=\frac1{|H|}\sum_{i=1}^m a_i\overline h_i
\]
on $N/H$ satisfies, by
\eqref{eq:p061-local-rpl-decomposition-candidate} and the $H$-invariance of
$u$,
\[
 q_N^*v
 =\frac1{|H|}\sum_{\gamma\in H}\gamma^*u
 =u|_N.
\]
The quotient $N^\circ/H$ is an ordinary open subset of both $N/H$ and
$C/H$.  Under \eqref{eq:p061-clean-local-quotient-candidate}, it is an
ordinary open neighbourhood of $x$.  Surjectivity of $q_N$ shows that $v$
restricts there to $g$.  Thus $g$ is $\mathbf R$-PL on an ordinary open
neighbourhood of every $x\in X$, and hence is $\mathbf R$-PL on $X$.
\end{proof}
